\section{What exactly is 
``AI sycophancy''?}
The term ``sycophancy'' has long-standing origins, and is thought to originate from the Ancient Greek term, ``tale-teller about figs'' \cite{Theoddes64:online}. This term has come to refer to excessive people-pleasing behavior of various forms, often for personal gain; Merriam Webster defines sycophancy as ``obsequious flattery''\footnote{\url https://www.merriam-webster.com/dictionary/sycophancy} and a sycophant as ``one who praises those in power in order to gain their approval'', while the Cambridge Dictionary defines sycophancy as ``behavior in which someone praises powerful or rich people in a way that is not sincere, usually in order to get some advantage from them'' \footnote{\url https://dictionary.cambridge.org/us/dictionary/english/sycophancy}. 

In practice, ``sycophancy'' is typically used in a negative manner to characterize those who aim to deceive or ingratiate themselves for personal gain. This presupposes some notion of intent. On the contrary, AI sycophancy research primarily takes a behaviorist standpoint, quantifying this phenomenon as excessive people-pleasing in AI systems that negatively impacts their outputs. Different works conceptualize ``excessive'' in different ways, and many works focus on accuracy or factuality to quantify the negative impacts of people-pleasing behavior [CITE]. Without a notion of ground truth, drawing this distinction becomes less trivial. A few works utilize normative frameworks including Bayesian rationality \cite{atwell2025basil}, while other works characterize ``excessive'' people-pleasing behavior as deviations from existing evidence: for instance, overturning negative paper reviews in response to user rebuttals, even when unsubstantiated \cite{wang2026ai}. Still other works define sycophancy as people-pleasing behaviors that lead to inconsistencies in LLMs' stated positions or outputs. 

This divergence between definitions reflects a fundamental disagreement among researchers regarding the types of model behaviors that constitute sycophancy. \citet{ye2026countsaisycophancytaxonomy} recently conducted a large-scale literature review and surveyed over 100 experts on AI sycophancy, and observed substantial variation in how sycophancy is defined and measured by researchers. The authors differentiate between position sycophancy, where models may shift their stated positions to align with or flatter users, and person sycophancy, where models may praise or positively evaluate users' traits and validate their emotions. They find the latter to be underrepresented in AI sycophancy research, with a lack of consensus among researchers regarding whether implicitly expressed Person behaviors should be considered sycophantic. A few works detail specific instances in which the lines may be blurred between optimizing for user preferences and sycophancy in AI systems. In particular, the line between personalization and sycophancy has been debated in the context of political discourse \cite{batzner2025germanpartiesqa}, as has the distinction between empathy and sycophancy \cite{PDFAIEmp14:online}. 

\section{AI sycophancy as norm displacement} One thing that many of these divergent definitions have in common, beyond the characterization of sycophancy as people-pleasing behavior, is a negative view of these behaviors or downstream impacts. Sometimes this is implicit, with references to AI sycophancy as an alignment failure \cite{Davalas2026MitigatingSA} or use of the term ``excessive'' in their characterizations of AI sycophancy \cite{Lee2026EffectsOA}. Sometimes these definitions are attached to certain negative impacts; for instance, many papers specifically define AI sycophancy as the tendency to conform to user opinions at the cost of accuracy 
\cite{Ibrahim2026TrainingLM} or in contrast with existing evidence \cite{Yao2026HearingIB, wang2026ai, Shah2026SycoPhantasyQS}. Other negative impacts mentioned in existing definitions of sycophancy include ``reinforcing biases'' \cite{lahlou2026writing}, ``providing false objectivity to transient anxieties'' \cite{PDFAIEmp14:online}, and encouraging behavior or beliefs that are ``aligned with a user's short-term preferences'' but ``detrimental to their long-term interests'' \cite{fulay2026delegates}. Others define sycophancy more broadly as the tendency to prioritize user approval above ``truth'' \cite{sharma2024towards, turner2026programmed}, which may generalize beyond notions of ground truth; \citet{turner2026programmed}, for instance, characterizes ``truth'' as ``an orientation toward 'getting things right''', not limited to specific truth values. Some researchers implicitly distinguish sycophancy as a negative side-effect of optimizing for user preferences; in other words, for outputs to be sycophantic, there must be a negative impact on output quality or negative downstream impacts that could emerge as a result of demonstrated people pleasing behaviors. For instance, \citet{PDFAIEmp14:online} argue that what makes this behavior sycophantic is the negative impact on certain users, specifically younger users.

Traditional notions of sycophancy in human behavior are fundamentally incompatible with the behaviorist lens through which AI sycophancy is viewed, and these two notions cannot be reconciled without, to some degree, anthropomorphizing AI systems. As a result, AI sycophancy researchers fundamentally disagree on how this concept should be defined and operationalized \cite{ye2026countsaisycophancytaxonomy}, with different works characterizing this phenomenon in wildly different ways. We argue that these divergent definitions can be reconciled by characterizing AI sycophancy broadly as \textit{model behaviors that accommodate a user's perceived belief, affect, identity, goal, or framing, in order to preserve approval or relational ease, while reducing fidelity to a context relevant norm}. In other words, rather than use the displacement of specific norms to characterize AI sycophancy (factuality, ``truth'', fairness, etc.), we define it as a general norm violation caused by apparent approval-seeking behavior in AI systems. This allows us to empirically verify which behaviors should fall under the umbrella of AI sycophancy, based on whether they reduce fidelity to context-dependent norms. This definition also accommodates significant variation regarding the types of behaviors researchers consider ``sycophantic'', while addressing growing concerns among researchers regarding one-size-fits-all approaches for labeling and addressing AI sycophancy \cite{Vennemeyer2025SycophancyIN, ye2026countsaisycophancytaxonomy}.



% \begin{enumerate}
%     \item 
%     \item Sycophancy research in AI
%     \begin{enumerate}
%         \item Origin of the popularization of ``AI sycophancy'', and metrics showing the trend in the number of sycophancy mentions
%         \item How, and why, AI sycophancy diverges from the dictionary definition
%         \item Overview of AI sycophancy research and proposed taxonomy (cite Meryl's paper)
%         \item Studied domains
%         \item Studied modalities
%         \item Studied mitigation strategies
%         \item Evaluation metrics
%     \end{enumerate}
%     \item AI sycophancy research in other domains
%     \begin{enumerate}
%         \item Psychology
%         \item Philosophy
%     \end{enumerate}
% \end{enumerate}